\begin{cvsection}{Scientific Projects}
    \cvitem
    {\href{https://gitlab.com/rhyme-org/ml_riemann_solver}{DeepRS}: A Feedforward Model for the Solution of the Riemann Problem}
    [Jan 2021 --- Present]
    [\ipa]
    % [In a typical Gudonov-based hydrodynamics simulation, the Riemann solver executes billions of times. Therefore, any improvement in this solver would result in a significant overall improvement. Using DeepRS, I so far managed to reach the accuracy of an exact solver with a substantial performance boost.]
    
    \cvitem
    {\href{https://gitlab.com/rhyme-org/rhyme}{Rhyme}, A Grid-based Radiation-Hydrodynamics Code for Astrophysics}
    [Mar 2018 --- July 2020]
    [\ipa]
    
    \cvitem
    {\href{https://github.com/saeedSarpas/HaloMatcher}{Cosmological Simulation Halo Matcher}}
    [May 2016 --- May 2017]
    [\AIfA]
    [Supervisor: Prof.\ Dr.\ Cristiano Porciani]
    % [A tool for retrieving \textit{matched} halos at different resolutions of cosmological simulations ran based on similar initial condition.]
    
    \cvitem
    {\href{https://gitlab.com/saeedSarpas/MatterPSGenerator}{A C Code to the Matter Power Spectrum}}
    [Nov 2015 --- Feb 2016]
    [\AIfA]
    [Supervisor: Prof.\ Dr.\ Cristiano Porciani]
    % [A modular C code for calculating the matter power spectrum of a given density field, extendable to calculate bi-spectrum.]
    
    \cvitem
    {\href{https://bitbucket.org/saeedSarpas/nbody-simulation/src/master}{Simulation of Gravitational N-body Systems Fully Implemented on Graphics Processing Units with CUDA/C Technology}}
    [May 2012]
    [\IASBS]
    % [Presented in the 16th annual meeting on research in Astronomy]
\end{cvsection}